
\setcounter{footnote}{0}
\chapter*{Expos\'e V\\Construction of Quotient Schemes}
\phantomsection\label{exp:sga3-V}
\addcontentsline{toc}{chapter}{Expos\'e V. Construction of Quotient Schemes}
\markboth{EXPOS\'E V. CONSTRUCTION OF QUOTIENT SCHEMES}%
  {EXPOS\'E V. CONSTRUCTION OF QUOTIENT SCHEMES}

\begin{center}
  \textit{by P. Gabriel}
\end{center}

The purpose of this expos\'e is to prove the theorems stated in TDTE~III.%
\footnote{\label{note:sga3-V-1}Editor's note: namely, \SGARefLink{rem:sga3-V-5.1}{Theorems~5.1}, 5.3,
\SGARefLink{lem:sga3-V-6.1}{6.1}, 6.2, and~\SGARefLink{lem:sga3-V-7.2}{7.2} of TDTE~III. The first two (respectively, the next two)
correspond to \hyperref[thm:sga3-V-4.1]{Theorem~4.1} (respectively, \hyperref[thm:sga3-V-7.1]{Theorems~7.1} and~\hyperref[thm:sga3-V-8.1]{8.1}) of this
expos\'e. \SGARefLink{lem:sga3-V-7.2}{Theorem~7.2} of TDTE~III is proved in Expos\'e~VI\textsubscript{A},
3.2 and~3.3.}
If \(X\) and \(T\) are two objects of a category~\(\mathscr C\), we write
\(X(T)\) instead of \(\operatorname{Hom}_{\mathscr C}(T,X)\). Similarly, if
\(\varphi:Y\to X\) is an arrow (respectively, if \(T\) is an object) of
\(\mathscr C\), then \(\varphi(T)\) denotes the map
\(g\mapsto\varphi\circ g\) from \(Y(T)\) to \(X(T)\):

\SGAThreeDiagram[0.23\textwidth]
  {figures/exp5/p251_phi_of_T_triangle.png}
  {The map \(\varphi(T):Y(T)\to X(T)\).}
  {fig:sga3-exp5-p251-phi-of-T}
  {Polo--Gille Exp5-13oct24.pdf, component-local p. 1; combined-reader p. 263; printed p. 251}

Likewise, \(T(\varphi)\) denotes the map \(g\mapsto g\circ\varphi\) from
\(T(X)\) to \(T(Y)\):

\SGAThreeDiagram[0.23\textwidth]
  {figures/exp5/p251_T_of_phi_triangle.png}
  {The map \(T(\varphi):T(X)\to T(Y)\).}
  {fig:sga3-exp5-p251-T-of-phi}
  {Polo--Gille Exp5-13oct24.pdf, component-local p. 1; combined-reader p. 263; printed p. 251}

Finally, if \(P\) is a scheme, we denote by \(\underline P\) the underlying
set of~\(P\).

As an exception, in the present expos\'e we do not follow the convention on
quotient notation stated in \SGARefLink{sga3:iv:scholium:4:6:15}{IV~4.6.15} (\emph{loc. cit.}, top of p.~227 of
the original), because we wish to give here a construction of quotients that
also applies to ``pre-equivalence relations''%
\footnote{\label{note:sga3-V-2}Editor's note: that is, groupoids with base
\(X\); see the terminology at the end of \hyperref[sec:sga3-V-1]{Section~1}. When \(\mathscr C\) is
the category of schemes, the quotient \(p:X\to Y\) of a groupoid \(X_\ast\)
with base \(X\) exists under certain hypotheses (see \hyperref[thm:sga3-V-4.1]{4.1}, \hyperref[lem:sga3-V-6.1]{6.1}, and~\hyperref[thm:sga3-V-7.1]{7.1}). If,
in addition, \(X_\ast\) is an equivalence relation, then under the same
hypotheses \(p\) is faithfully flat and quasi-compact, hence a universal
epimorphism; see \emph{loc. cit.}.}
that are not equivalence relations.

\section*{1. \(\mathscr C\)-groupoids}
\phantomsection\label{sec:sga3-V-1}
\addcontentsline{toc}{section}{1. \texorpdfstring{\(\mathscr C\)}{C}-groupoids}

\phantomsection\label{par:sga3-V-1-a}%
\noindent\textbf{a)} Let \(\mathscr C\) be a category in which products and
fiber products exist. Recall first that a diagram

\SGAThreeDiagram[0.39\textwidth]
  {figures/exp5/p252_exact_pair_cokernel.png}
  {A parallel pair followed by its proposed cokernel.}
  {fig:sga3-exp5-p252-exact-pair-cokernel}
  {Polo--Gille Exp5-13oct24.pdf, component-local p. 2; combined-reader p. 264; printed p. 252}

\noindent in \(\mathscr C\) is said to be \emph{exact} if
\(p d_0=p d_1\) and if, for every \(T\in\mathscr C\), the map \(T(p)\) is a
bijection from \(T(Y)\) onto the subset of \(T(X_0)\) consisting of the
arrows \(f:X_0\to T\) such that \(f d_0=f d_1\). One also says that
\((Y,p)\) is the \emph{cokernel} of \((d_0,d_1)\), and writes
\phantomsection\label{eq:sga3-V-1-cokernel}%
\begin{equation*}
  (Y,p)=\operatorname{Coker}(d_0,d_1).
\end{equation*}

\phantomsection\label{par:sga3-V-1-b}%
\noindent\textbf{b)} For example, let \(\mathscr C\) be the category
\((\mathbf{Esp.An})\) of ringed spaces. In this case a cokernel \((Y,p)\)
always exists, and it may be described as follows. The underlying topological
space of \(Y\) is obtained from \(X_0\) by identifying the points \(d_0(x)\)
and \(d_1(x)\), and by giving \(Y\) the quotient topology. The canonical map
\(\pi:X_0\to Y\), together with \(d_0,d_1\), then induces a pair of arrows
of sheaves of rings on \(Y\):

\SGAThreeDiagram[0.66\textwidth]
  {figures/exp5/p252_ringed_space_sheaf_pair.png}
  {The pair of induced arrows between direct-image sheaves.}
  {fig:sga3-exp5-p252-ringed-space-sheaf-pair}
  {Polo--Gille Exp5-13oct24.pdf, component-local p. 2; combined-reader p. 264; printed p. 252}

Here \(\mathcal O_i\) is the structure sheaf of \(X_i\). As the sheaf of
rings on \(Y\), one chooses the subsheaf of \(\pi_*(\mathcal O_0)\) whose
sections \(s\) satisfy \(\delta_0(s)=\delta_1(s)\). The arrow \(p\) is
defined in the evident way.

Consider%
\footnote{\label{note:sga3-V-3}Editor's note: \hyperref[lem:sga3-V-1.1]{Lemmas~1.1} and~\hyperref[sga3:exp5:lemma:1.2]{1.2} have been
added; they are used several times in Sections~\hyperref[sec:sga3-V-5]{5}--\hyperref[sec:sga3-V-9]{9}.}
the following diagram in \((\mathbf{Esp.An})\):

\SGAThreeDiagram[0.18\textwidth]
  {figures/exp5/p252_ringed_space_pair_inline.png}
  {The pair \((d_0,d_1):X_1\rightrightarrows X_0\).}
  {fig:sga3-exp5-p252-ringed-space-pair}
  {Polo--Gille Exp5-13oct24.pdf, component-local p. 2; combined-reader p. 264; printed p. 252}

\noindent Let \((Y,p)\) be its cokernel. An open subset \(U\) of \(X_0\)
is said to be \emph{saturated} if
\[
  d_0^{-1}(U)=d_1^{-1}(U),
\]
which is equivalent to saying \(U=p^{-1}(p(U))\). In this case, since \(Y\)
has the quotient topology, \(p(U)\) is an open subset of~\(Y\).

\medskip
\phantomsection\label{lem:sga3-V-1.1}%
\noindent\textbf{Lemma 1.1.}
\textit{Let \(U\) be a saturated open subset of \(X\), and set
\(V=p(U)\). Denote by
\(U_1=d_0^{-1}(U)=d_1^{-1}(U)\) the corresponding open subset of \(X_1\),
and by \(\widetilde d_0,\widetilde d_1\), and \(\widetilde p\) the
restrictions of \(d_0,d_1\) to \(U_1\), and of \(p\) to \(U\), respectively.
Then \((V,\widetilde p)\) is a cokernel in \((\mathbf{Esp.An})\) of:}%
\footnote{\label{note:sga3-V-4}Editor's note: this is not the case in the
category of schemes. For example, let \(S=\operatorname{Spec}(\mathbf C)\),
let
\(X_0=\mathbf A^2_S=\operatorname{Spec}(\mathbf C[x_1,x_2])\), let
\(d_1:\mathbf G_{m,S}\times_S\mathbf A^2_S\to\mathbf A^2_S\) be the
action of \(\mathbf G_{m,S}\) on \(\mathbf A^2_S\) by homotheties, let
\(d_0\) be the projection onto the second factor, and let
\(U=\mathbf A^2_S\setminus\{m\}\), where \(m=(0,0)\). Then the projective
line \(\mathbf P^1_S\) is the cokernel of
\((\widetilde d_0,\widetilde d_1)\) in both \((\mathbf{Esp.An})\) and
\((\mathbf{Sch})\), whereas the cokernel \(Y\) of \((d_0,d_1)\) in
\((\mathbf{Esp.An})\) is the union of \(\mathbf P^1_S\) and the point
\(y_0=\{p(m)\}\). The only open subset containing \(y_0\) is \(Y\), and
\(\Gamma(Y,\mathcal O_Y)=\mathbf C\). If
\(f:\mathbf A^2_S\to T\) is a morphism of \(S\)-schemes such that
\(f d_0=f d_1\), and if \(V=\operatorname{Spec}(A)\) is an affine open
subset of \(T\) containing \(t_0=f(y_0)\), then
\(f^{-1}(V)=\mathbf A^2\), and the ring homomorphism
\(A\to\mathbf C[x_1,x_2]\) factors through \(\mathbf C\). This shows that
\(S=\operatorname{Spec}(\mathbf C)\) is the cokernel of \((d_0,d_1)\) in
the category \((\mathbf{Sch}/S)\).}

\SGAThreeDiagram[0.40\textwidth]
  {figures/exp5/p252_lemma_1_1_restricted_cokernel.png}
  {The restricted pair and its cokernel in \hyperref[lem:sga3-V-1.1]{Lemma~1.1}.}
  {fig:sga3-exp5-p252-lemma-1-1-restricted-cokernel}
  {Polo--Gille Exp5-13oct24.pdf, component-local p. 2; combined-reader p. 264; printed p. 252}

The verification is easy.
